% Sean P. Florez | curriculum vitae
% Layout and typography follow Sean's supplied LaTeX template.
% Compile: pdflatex main.tex
\documentclass[letterpaper,10pt]{article}
\usepackage[left=0.65in,right=0.65in,top=0.55in,bottom=0.75in]{geometry}
\usepackage{enumitem}
\usepackage[colorlinks=true,urlcolor=blue]{hyperref}
\usepackage{titlesec}
\usepackage{parskip}
\usepackage{amsmath}
\usepackage{needspace}

% Searchable text; the default Computer Modern typography is unchanged.
\input{glyphtounicode}
\pdfgentounicode=1
\hypersetup{
    pdftitle={Sean P. Florez | Curriculum Vitae},
    pdfauthor={Sean P. Florez}
}

% Section spacing and rule, as supplied.
\setlength{\parindent}{0pt}
\setlength{\parskip}{3pt}
\titlespacing\section{0pt}{10pt plus 2pt minus 1pt}{5pt plus 1pt minus 1pt}
\titleformat{\section}{\large\bfseries}{\thesection}{1em}{}[\titlerule]

% List style, as supplied.
\setlist[itemize]{leftmargin=*,noitemsep,topsep=0pt}
\setlength{\emergencystretch}{1em}
\widowpenalty=10000
\clubpenalty=10000
\raggedbottom
\hyphenation{Quantum materials manufacturing}

% Page-break guards: keep a heading with the lines that follow it.
\newcommand{\keep}{\needspace{4\baselineskip}}
\newcommand{\keepsection}{\needspace{6\baselineskip}}

\begin{document}
\pagestyle{plain}

% ===== Header =====
\begin{center}
    {\LARGE \textbf{Sean P.\ Florez}}\\
    \href{mailto:sean.florez@colorado.edu}{sean.florez@colorado.edu} \,\textbar\, (786) 202--5558 \,\textbar\, Boulder, CO\\
    \href{https://seanflorez.com}{seanflorez.com} \,\textbar\,
    \href{https://linkedin.com/in/sean-florez}{linkedin.com/in/sean-florez} \,\textbar\,
    \href{https://github.com/fl-sean03}{github.com/fl-sean03}
\end{center}

% ===== Education =====
\section*{Education}
\textbf{University of Colorado Boulder} \hfill \textit{Expected May 2028}\\
Ph.D., Materials Science \& Engineering. Advised by Prof.\ Hendrik Heinz in the Heinz Interfaces Laboratory.\\
Atomistic simulation of interfaces; experiment-calibrated force fields; autonomous-laboratory infrastructure.

\textbf{University of Florida} \hfill \textit{May 2024}\\
B.S., Materials Science \& Engineering. Minor in Chemistry.

% ===== Research =====
\keepsection\section*{Research Experience}
\textbf{Graduate Research Assistant} \hfill \textit{University of Colorado Boulder, Oct 2024--Present}\\
\textit{Heinz Interfaces Laboratory, Prof.\ Hendrik Heinz}

\keep\textit{Interlayer shear in MXenes (lead author, manuscript in preparation)}
\begin{itemize}
    \item Ran the production shear campaign for OH-, F-, and mixed-terminated Ti\(_3\)C\(_2\) bilayers in LAMMPS with \textbf{780 NVT trajectories} across six termination and geometry cells on CU's Alpine cluster, ten seeds per stress, zero failed runs.
    \item Froze the protocol before production and had the campaign audited read-only afterward. The audit re-derived every reported critical stress from raw trajectories using an independent re-implementation of the slip classifier, and caught an inflated trajectory count in our own reporting before it reached the manuscript.
    \item Termination sets interlayer shear resistance across roughly a threefold range (\(\tau_{0.5}\) 42--135 MPa), OH \(>\) mixed \(>\) F.
\end{itemize}

\keep\textit{Hydrogen release from N-ethylcarbazole on Pt nanocrystals (lead author, manuscript in preparation)}
\begin{itemize}
    \item Designed and ran the MD campaign behind the manuscript, \textbf{60 slab production runs} across three Pt facets plus cuboctahedral nanoparticle ensembles, zero production job failures.
    \item Found that dehydrogenation-reactive configurations concentrate at low-coordination sites. Vertices carry 5.4\(\times\) the per-atom contact density of \{100\} terraces and 2.9\(\times\) that of \{111\}, so local coordination, rather than facet identity, sets where release happens.
    \item Built the per-area normalization methodology that makes facets, edges, and vertices comparable, and validated the finite nanoparticle against infinite slabs facet by facet.
    \item Wrote the technical narrative for the ALCF Director's Discretionary allocation that funds the next campaign.
\end{itemize}

\keep\textit{Copper-oxide force fields for INTERFACE FF (industry-funded program)}
\begin{itemize}
    \item Parameterized INTERFACE force-field models for copper oxides as the first-year university deliverable, calibrated exclusively against experimental lattice constants, elastic constants, surface energies, and vibrational spectra.
    \item Delivered a validated Cu\(_2\)O parameter family in bonded 9-6 and 12-6 forms, anchored to experimental lattice constants and within \textbf{\(\pm\)5\%} of both independent experimental elastic references. Shipped as a distribution bundle (LAMMPS data files and run decks, an \texttt{.frc} for msi2lmp, a CHARMM/NAMD \texttt{.prm}), with all \textbf{52 configurations} run-verified.
\end{itemize}

\keep\textit{Retrieval over the force-field corpus}
\begin{itemize}
    \item Built a SOAP-embedding retrieval index over the group's \textbf{3,009-structure, 6.7M-atom} corpus so a new atom's type and parameters can be proposed from its nearest structural neighbors behind a confidence gate. The same index maps coverage gaps and audits type consistency.
\end{itemize}

\keep\textbf{High-Performance Computing Intern} \hfill \textit{Air Force Research Laboratory, May--Aug 2025}\\
\textit{Materials \& Manufacturing Directorate, Wright-Patterson AFB}
\begin{itemize}
    \item Built the LAMMPS shear-testing workflow for OH-terminated Ti\(_3\)C\(_2\) bilayers under controlled hydration and ran the stress sweeps on DoD HPC.
    \item Fit slip probability against applied shear stress with logistic models. Hydration turns out to be non-monotonic. Dry bilayers hold to 103 MPa, a quarter monolayer of water drops them to 8 MPa, and full coverage recovers to 39 MPa as water rebuilds bridging hydrogen bonds.
    \item First MD evidence that hydration tunes MXene interlayer shear strength by an order of magnitude, with the stress windows that follow for EM-shielding composite design.
\end{itemize}

\keep\textbf{Undergraduate Research Assistant} \hfill \textit{Hennig Group, University of Florida, Oct 2021--May 2024}
\begin{itemize}
    \item Benchmarked the VASPsol implicit solvation model against experimental energies from the Minnesota Solvation Database for neutral and charged solutes in water.
    \item Optimized cavity and screening parameters by grid search and Nelder--Mead across a large DFT campaign, and outlined a surrogate-model strategy that replaces high-count parameter sweeps for future calibrations.
\end{itemize}

\keep\textbf{Materials Modeling Intern} \hfill \textit{Air Force Research Laboratory, May--Dec 2023}
\begin{itemize}
    \item Built and validated an automated DFT workflow (Quantum ESPRESSO, ASE, and pymatgen) for high-throughput screening of low-work-function emitters, giving the group a reproducible selection path for field-emission candidates.
\end{itemize}

% ===== Research automation =====
\keepsection\section*{Research Automation \& Compute Infrastructure}
\textbf{OpenSDL, computational and autonomous laboratory framework} \hfill \textit{2026--Present}\\
\href{https://github.com/fl-sean03/OpenSDL}{github.com/fl-sean03/OpenSDL} \;\textbar\; \href{https://seanflorez.com/OpenSDL/}{seanflorez.com/OpenSDL} (documentation) \;\textbar\; Apache-2.0
\begin{itemize}
    \item An open foundation for building computational and autonomous laboratories. Declare what a laboratory can do, execute those declarations as reproducible workflows, preserve the evidence, and feed the result back into the choice of the next experiment. Roughly 26,000 non-test lines of Python, about 600 test functions, and CI.
    \item A durable runtime with DAG execution, retries, timeouts, resource leases, restart reconciliation, and policy checks, over SQLite metadata and content-addressed artifact storage. Adapters for a simulated lab, local compute, and human tasks, with Slurm and hardware integration in progress; a CLI, Python SDK, HTTP API, and MCP hook; portable run export.
    \item Implemented campaign scoring and next-round selection. A reference three-dye formulation search matched an unseen target recipe within \textbf{0.5 percentage points}, evaluating \textbf{96 candidates per round}.
\end{itemize}

\keep\textbf{agent-fleet, persistent research automation} \hfill \textit{2026--Present}\\
\href{https://github.com/fl-sean03/agent-fleet}{github.com/fl-sean03/agent-fleet}
\begin{itemize}
    \item Built a research harness based on a \textbf{15-agent fleet} operating for months, with isolated workspaces, persistent conversations, and inter-agent messaging.
    \item Ran multi-day computational campaigns with recovery after reboots, credential expiry, and preemption.
\end{itemize}

\keep\textbf{HPC \& Cloud Throughput Infrastructure} \hfill \textit{2025--Present}\\
\href{https://github.com/fl-sean03/subjob}{github.com/fl-sean03/subjob} \;\textbar\; \href{https://github.com/fl-sean03/cloudcomputemanager}{github.com/fl-sean03/cloudcomputemanager}
\begin{itemize}
    \item \textbf{subjob} runs many heterogeneous tasks inside one HPC allocation, so a hundred short jobs cost one queue wait instead of a hundred. Shared task pool, workers spanning nodes and partitions, per-task walltime enforcement, and heartbeat-driven recovery of a dead worker's claims. Built for the Pt campaign's per-snapshot analysis fan-outs.
    \item \textbf{CloudComputeManager} (MIT) manages GPU-cloud lifecycles for scientific workloads on spot instances, covering provisioning, environment setup, multi-stage pipelines, checkpoint and preemption recovery, batch sweeps, and cost tracking. Workload-agnostic, with 368 tests. Used for LAMMPS shear campaigns.
\end{itemize}

% ===== Manufacturing =====
\keepsection\section*{Manufacturing \& Production Systems}
\textbf{Engineering Contractor, Rocky Mountain Scientific Laboratory} \hfill \textit{Jun--Aug 2026}\\
\textit{Quality system for a client affiliate's new energetic-materials production line}\\
\textit{Automation \& Robotics embed, Jun--Aug 2026}
\begin{itemize}
    \item Developed lot-release records covering inspections, approval authority, hold criteria, and nonconformance disposition, with item-level traceability to the governing \textbf{military specification}.
    \item Built a controlled, single-source document set with \textbf{10 client-format instructions}, \textbf{20 test-specific quality-control packets}, manufacturing procedures, and inspections.
    \item Automated checks for process-route isolation, link validity, terminology, and consistency between released documents and their source.
    \item Built a program tracker that generates the Kanban board, decision log, requirements traceability matrix, and design-review record.
    \item Supported quality, packaging, labeling, on-site commissioning, and hazard and failure-mode analysis.
\end{itemize}

% ===== National security =====
\keepsection\section*{National Security \& Technology Transition}
\textbf{S\&T Scouting Intern} \hfill \textit{The Joint Staff, J7 / Future Technology Office, Aug--Oct 2025}
\begin{itemize}
    \item Prioritized propulsion and energy concepts and produced \textbf{eight technical briefs} assessing readiness, risk, and experimentation paths for TRL 3--5 technologies.
    \item Coordinated with Service laboratories and DoD R\&D organizations to connect emerging academic research to operational experimentation and transition pipelines.
\end{itemize}

\keep\textbf{Technology Commercialization Intern} \hfill \textit{Idaho National Laboratory (DOE), May--Aug 2024}
\begin{itemize}
    \item Ran \textbf{20 customer-discovery interviews} under Energy I-Corps and authored the commercialization plan for a laboratory-origin technology.
    \item Built AI and natural-language tooling for the technology-transfer team, and supported a DOE proposal that was awarded \textbf{\$2M}.
\end{itemize}

% ===== Leadership =====
\keepsection\section*{Leadership \& Service}
\textbf{Co-Founder, \href{https://www.lablinkinitiative.org}{LabLink Initiative}} \hfill \textit{Aug 2024--Present}
\begin{itemize}
    \item Nonprofit moving community college students into national-laboratory research internships, principally the DOE Office of Science CCI and SULI programs. Partnerships with three Sacramento-area community colleges, and a webinar delivered across Phi Theta Kappa, the honor society for two-year colleges, that drew more than \textbf{250 students}.
    \item Built the matching pipeline and the tooling behind it.
\end{itemize}

% ===== Publications and outputs =====
\keepsection\section*{Publications \& Outputs}
\textbf{Co-author:} \textit{INTERFACE Force Field for Alumina with Validated Bulk Phases and a pH-Resolved Surface Model Database for Electrolyte and Organic Interfaces.} \href{https://arxiv.org/abs/2601.12570}{arXiv:2601.12570}, 2026 (preprint).

\keep\textbf{Lead-author manuscripts in preparation:}
\begin{itemize}
    \item \textit{Modeling of MXene Shear Properties via Atomistic Molecular Dynamics Simulations}, with AFRL Materials \& Manufacturing Directorate co-authors.
    \item \textit{Understanding Hydrogen Release Mechanisms of N-Ethylcarbazole on Pt Nano-Crystal Catalysts}.
\end{itemize}

\textbf{Open-source releases:} \href{https://github.com/fl-sean03/MolSAICV4}{MolSAIC}, a code-first MD system builder for atomistic model construction and force-field assignment; \href{https://github.com/fl-sean03/pdb2msi}{pdb2msi}, a PDB to Materials Studio structure converter.

\textbf{Computing allocations:} Argonne Leadership Computing Facility Director's Discretionary award \textit{HydrogenStorage}, approximately 26,150 node-hours across Polaris, Crux, Sophia, and Graphcore, awarded May 2026; CU Boulder Alpine and DoD HPC allocations for production MD campaigns, including 480-job Slurm arrays.

\keep\textbf{Quantum-chemical property prediction on QM9} \hfill \href{https://github.com/fl-sean03/qm9-smiles-predictor}{GitHub: qm9-smiles-predictor}
\begin{itemize}
    \item Four reproducible pipelines (single-task, multitask, autoencoder, hybrid Mordred \(+\) ECFP4) predicting nine quantum-chemical properties across roughly 130k molecules.
    \item Reproduced and critiqued published SMILES-FNN baselines, showing where multitask learning underperforms and where hybrid descriptors reduce MAE over Mordred-only models, especially for dipole moment and polarizability.
\end{itemize}

\keep\textbf{Automated phase identification in experimental powder XRD} \hfill \href{https://github.com/fl-sean03/opxrd-ml-binary-phase}{GitHub: opxrd-ml-binary-phase}
\begin{itemize}
    \item Ingested the roughly 92k-pattern opXRD archive (about 2.4\% labeled) and formulated binary classification for a rare phase, with an interpolated 2\(\theta\) grid and normalized intensities.
    \item Trained logistic-regression and random-forest classifiers with SMOTE for extreme class imbalance, reaching a 0.80 minority-phase F1 with balanced precision and recall on held-out data. A prototype for automated pXRD analysis in self-driving laboratories.
\end{itemize}

% ===== Skills =====
\keepsection\section*{Technical Skills}
\textbf{Molecular dynamics:} LAMMPS, NAMD, Materials Studio/Discover. Model building, equilibration protocol design, shear and deformation workflows, trajectory analysis at multi-GB scale.\\
\textbf{Electronic structure:} VASP, VASPsol, Quantum ESPRESSO. High-throughput workflows via ASE and pymatgen.\\
\textbf{Force fields:} INTERFACE FF, CVFF, CHARMM, PCFF, COMB3, EAM. Parameterization against experimental data, transferability testing, elastic and surface-energy validation.\\
\textbf{Machine learning:} scikit-learn, TensorFlow, imbalanced-learn, RDKit, Mordred, SOAP descriptors. Regression, classification, class-imbalance handling, surrogate modeling, uncertainty and error analysis.\\
\textbf{Python:} NumPy, SciPy, pandas, MDAnalysis, matplotlib. Package design, CLI tooling, pytest, ruff, pyright.\\
\textbf{HPC and systems:} Slurm and PBS Pro, job arrays, checkpoint-restart, throughput benchmarking, GPU and CPU-MPI builds; Linux, bash, git, systemd, SQLite, Docker, GitHub Actions, MCP.\\
\textbf{Quality and manufacturing:} Lot-release record and human-factors document design, military-specification traceability, requirements traceability matrices, PDR/CDR stage gates, FMEA and process hazard analysis, BOM development, vendor qualification, labeling compliance.

% ===== Clearance =====
\keepsection\section*{Security Clearance}
Top Secret (active); SCI eligible.

\end{document}
